slavabogy
\documentclass[10pt]{article}

\def\Type{LessonDJ}
\def\TypeNum{Workshop Week 13} %Lesson #
\def\TypeTitle{FTOC and Intro to Probability} %Lesson Title
\def\Semester{Fall 2021} %Not Used 
\def\Course{MATH 1190}%Course number and name
\def\Targets{    \target{I2}, \target{I3}, \target{F3} }

\usepackage{preambleDJ} 

%%%%%%%%%%%% BEGIN DOCUMENT %%%%%%%%%%%%%%%
%%%%%%%%%%%% BEGIN DOCUMENT %%%%%%%%%%%%%%%
%%%%%%%%%%%% BEGIN DOCUMENT %%%%%%%%%%%%%%%
%%%%%%%%%%%% BEGIN DOCUMENT %%%%%%%%%%%%%%%
%%%%%%%%%%%% BEGIN DOCUMENT %%%%%%%%%%%%%%%

\begin{document}
\pagestyle{otherpages}
\thispagestyle{firstpage}


\vs{.6}

\section{Warm-Up}
\begin{enumerate}
    \item Find the following anti-derivatives 
    \begin{enumerate}
        \item $\ds \int x^3 \;dx$
        \item $\ds \int \frac{1}{2x} \;dx$
        \item $\ds  \int 4\cdot3^x\;dx$
    \end{enumerate}
    \item Compute the following definite integral
    \[\int_2^5 \frac{x^4+1 + x3^x}{2x}dx\]
    \item Compute the following definite integrals:
    \begin{enumerate}
        \item $\ds \int_0^3 (2x+1)^2\,dx.$
         \item $\ds \int_1^4 \left(\frac 1x + 2\right)\left(3x + 1\right)\,dx.$
    \end{enumerate}    

\end{enumerate}

\newpage
 \section{More Definite Integrals}
   \begin{enumerate}
    \item Consider the following integral:
    \[\int_{-1}^3\frac{1}{x^4}dx\]
    Does the fundamental theorem of calculus apply? If not, explain why not. If so, find the definite integral.\\
    \item Consider the function $v(t) = 4t^3-16t^2+16t$ shown below:
    \begin{center}
        \begin{tikzpicture}
          \begin{axis}[
            xmin=-2.5,        % Minimum x-value
            xmax=2.5,         % Maximum x-value
            ymin=-8,         % Minimum y-value
            ymax=8,          % Maximum y-value
            axis lines=middle, % Draw the x and y axes in the middle
            x label style={anchor=west},
            xlabel={$t$},     % Label for the x-axis
            ylabel={$v(t)$},     % Label for the y-axis
            xtick       = {-2,-1,1,2},
            ytick       = {-8,-4,4,8},
            %yticklabels = {$2$,$4$,$6$,$8$},
            axis line style={thick, <->}, % Add arrows at the end of the axes
            width = 2in, height=2in,
            %restrict y to domain=-6.5:6.5
          ]
          %\addplot[domain=-2.5:2.5, samples=400, blue, ultra thick] {x^3};
              \addplot[name path=f1, domain=-2.5:2.5, blue, ultra thick, <->] {4*x^3-16*x^2+16*x};
%  \path[name path=axis] (axis cs:-2.5,0) -- (axis cs:2.5,0);
%     \addplot [
%        thick,
%        color=blue,
%        fill=blue, 
%        fill opacity=0.1
%    ]
%    fill between[
%        of=f1 and axis,
%        soft clip={domain=-1:2},
%   ];
          \end{axis}
        \end{tikzpicture}
    \end{center}

    \begin{enumerate}
        \item[a)] What is the area bounded by $v$ and the $t$-axis from $t = 1$ to $t = 2$?
        \item[b)] What is the net area bounded by $v$ and the $t$-axis from $t = -1$ to $t = 1$?
        \item[c)] What is the net area bounded by $v$ and the $t$-axis from $t = -1$ to $t = 2$?
    \end{enumerate}

    \item Suppose $v$ represented the velocity of a car as a function of time. What do the quantities you calculated in the previous part represent?
    
\end{enumerate}

\newpage
\section{Probability}
\begin{enumerate}
\item Peter T, a pre-dentist student,  is a Lawn resident and was excited to participate in this year's Trick-or-Treating on the Lawn. He decided to offer Trick-or-Treaters his favorite candy bar, Snickers, as well as something, um, ``more practical"-- toothbrushes.   He put $2$ Snicker bars and $3$ toothbrushes into a grab bag and asked each Trick-or-Treater to pick two items without looking and without the ability to put any back. What's the chance a Trick-or-Treater  walked away crying because they got  $0$ Snickers (i.e $2$ toothbrushes?  What's the chance the Trick-or-Treater walked away jumping for joy because they got $2$ Snickers? What's the chance they walked away confused because they got $1$ Snickers (and thus $1$ toothbrush)?

Being a part of the famed Math $1190$ crew, you decide to create a histogram to figure out these probabilities.  
\begin{enumerate}
\item What is the sample space for pulling two items out of the bag?
\item You  let $X$ be a random variable that describes how many Snickers the Trick-or-Treater received.  What are the possible outcomes for $X$?
\item What is $P(X=0)$?  This is the chance the Trick-or-Treater gets $0$ Snickers (and thus  $2$ toothbrushes) and bursts out in tears. The next $3$ steps will help you find  $P(X=0)$.
\enum{
\item What's the chance the Trick-or-Treater will get a toothbrush on the first grab?
\item Now think about how many Snickers and how many toothbrushes remain in the bag after the first grab. What's the chance to get a toothbrush on the 2nd grab?
\item Multiply these two results together to find $P(X=0)$
}
\item What is $P(X=2)$? This is the probability the Trick-or-Treater will jump for joy because they received $2$ Snickers. The next $3$ steps will help you find  $P(X=2)$.
\enum{
\item What is the chance the Trick-or-Treater will get a Snickers on the first grab?
\item Now think about how many Snickers and how many toothbrushes remain in the bag after the first grab. What's the chance to get a Snickers on the 2nd grab?
\item Multiply these two results together to find $P(X=2 )$
}
\item What is $P(X=1)$? This is the chance that the Trick-or-Treater walked away confused because they got one Snickers and one toothbrush. {\em Hint:} There are at least 2 different approaches to come up with this probability.
\item Create a histogram that represents all of this information.\\
\begin{tikzpicture}
    \begin{axis}[
        ybar=0, % Specifies vertical bars
        ymin=0, % Ensures bars start from zero
        xtick=data, % Places x-axis ticks at data points
        xticklabels={,,}, 
                  %   xticklabels={0,1,2}, % Custom x-axis labels
                %   ytick       = {0.1,0.3,0.6},
            %yticklabels =  {$0.1$, $0.3$, $0.6$},
        enlarge x limits={0.2}, % Adjusts spacing between bars and axis limits
            y label style={anchor=south west},
            ylabel={$p(x)$},
                        xlabel = {$x$},
		x label style={anchor=west},
		 ymin = 0, ymax = .6,
        bar width=20pt, % Adjusts the width of the bars
                    width = 2.5in, height=1.4in,
       % legend style={at={(0.5,-0.15)},anchor=north,legend columns=-1} % Legend positioning
    ]
  \addplot[fill=blue!40] coordinates {(1,0)  (2,0)  (3,0) };
  %   \addplot[fill=blue!40] coordinates { (1,0.3)  (2,0.6)  (3,0.1) };
    \end{axis}
\end{tikzpicture}
\end{enumerate}


    \item Suppose two standard dice are rolled and the result is recorded.
\begin{enumerate}
    \item[a)] Let $X$ be the difference of first die minus the second die. 
    \enum{
    \item Construct a table of all possible combinations of rolls of dice and the resulting difference similar to the one we used in class.
    \item Find the probability for all possible outcomes for $X$
    \item Sketch a histogram
    }
    \end{enumerate}
    \mpageT{.5}{.5}[\hs{.2}]{
    \begin{tikzpicture}
    \begin{axis}[
        ybar, % Specifies vertical bars
        bar shift=0,
        ymin=0, % Ensures bars start from zero
        xtick=data, % Places x-axis ticks at data points
                     xticklabels={,,,,,,,,,}, % Custom x-axis labels
                     ytick=\empty,
%                    ytick       = {1/36, 2/36, 3/36, 4/36, 5/36, 6/36},
%            yticklabels =  {$1/36$,$2/36$,$3/36$, $4/36$,$5/36$,$6/36$},
        enlarge x limits={0.2}, % Adjusts spacing between bars and axis limits
            y label style={anchor=south west},
            ylabel={$p(x)$},
                        xlabel = {$x$},
		x label style={anchor=west},
		 ymin = 0, ymax = .17,
        bar width=17pt, % Adjusts the width of the bars
                    width = 4in, height=1.4in,
                    %enlarge x limits={1.2}
       % legend style={at={(0.5,-0.15)},anchor=north,legend columns=-1} % Legend positioning
    ]
 \addplot[fill=blue!40] coordinates {(1,0)  (2,0)  (3,0) (4,0) (5,0) (6,0) (7,0) (8,0) (9,0) (10,0) (11,0)};
 %\addplot[fill=blue!40] coordinates {(1,0.027)  (2,0.055)  (3,0.083) (4,0.11) (5,0.139) (6,0.167) (7,0.139) (8,0.11) (9,0.083) (10,0.055) (11,0.027)};
    \end{axis}
\end{tikzpicture}
    }{
    \begin{enumerate}
    \item[b)] Compute $P(X \geq 2)$.
    \item[c)] Compute $P(X \leq -2)$.
    \item[d)] Compute $P(-1\leq X \leq 1)$.
    \item[e)] Compute $P(X \text{ is even})$.
    \end{enumerate}
    }


%\item Suppose in your sock drawer, you have a pair of red socks, a pair of black socks, and a pair of brown socks. On your way to school, you randomly pull out single socks (without replacement), until you have picked two of the same color. 
%\begin{enumerate}
%    \item[a)] What is the sample space?
%    \item[b)] Let $X$ be the number of socks you have to pull. Write down the probability mass function for $X$. 
%\end{enumerate}


\end{enumerate}



\end{document}